% !TEX encoding = IsoLatin

%%%%%%%%%%%%%%%%%%%%%%%%%%%%%%%%%%%%%%%%%%%%%%%%%%%%%%%%%%%%%%%
%                                                             %
% Una funci�n invariante respecto a un                        %
% grupo de transformaciones                                   %
% by                                                          %
% Juan Luis Varona                                            %
%                                                             %
%%%%%%%%%%%%%%%%%%%%%%%%%%%%%%%%%%%%%%%%%%%%%%%%%%%%%%%%%%%%%%%

%	I have made a picture that shows that the function $f(x)$ is invariant 
%	with respect to the group of transformations $\{H_0(x),\dots,H_5(x)\}$, 
%	that is, $f(H_i(x)) = f(x)$, $i=0,\dots,5$.
%	I think that this is property is known, but I have not found a reference.


\documentclass[autocontact]{gaceta}

%---------------------
% La Gaceta se har\'a con esto activado, 
% pero no hace falta ahora si no quieres:
\usepackage[T1]{fontenc}
\usepackage{lmodern}

% Puedes usar esto (u otra codificaci\'on) si quieres:
%\usepackage[latin1]{inputenc}
\usepackage[spanish]{babel}
%\decimalpoint
%---------------------

%---------------------
% Esto ya lo carga el estilo de La Gaceta:
%\usepackage{amsmath, amsthm, amssymb}
%\usepackage{url} % <-- Para p\'aginas web o similar: \url{...}
%\usepackage{graphicx}
%---------------------

%---------------------
% Carga esto para gr\'aficos si lo necesitas:
% \usepackage{wrapfig}
%---------------------

%---------------------
\usepackage{tikz}
%---------------------

%---------------------
% <<< ESTO SE AJUSTAR\'A AL EDITAR CADA N\'UMERO DE LA GACETA >>>
%\provisional
\newcounter{pageforminiatura}
\setcounter{pageforminiatura}{344} % Siempre ser\'a p\'agina par
\setcounter{page}{1} % No se usa (dejarlo como 1)
\journame{La Gaceta de la RSME}
\yearofpublication{2016}
\volume{19}
\issuenumber{2}
%---------------------

%---------------------
% Deja esto as\'{\i}:
%\belongstopart{Art\'{\i}culos} 
%\belongstopart{Secciones} 
%\belongstopart{Noticias de la Sociedad} 
%\belongstopart{Actualidad} 
\nopart
%---------------------

%---------------------
% Deja esto as\'{\i}: 
\supertitle{\large Miniaturas matem\'aticas de \thejourname\\[1mm]
\hfill\footnotesize
Vol.~\thevolume\ (\theyearofpublication), 
N\'um.~\theissuenumber, P\'ag.~\thepageforminiatura\theprovisional}
%
\noeditor
%---------------------

%---------------------
\title{Una funci�n invariante\\respecto a un grupo de transformaciones}
%\author{} % Poner \author{} o no poner nada, da igual.
\contact{Juan Luis Varona, 
Departamento de Matem\'aticas y Computaci\'on, 
Universidad de La Rioja}{jvarona@unirioja.es}
{http://www.unirioja.es/cu/jvarona/}
%---------------------

%%%%%%%%%%%%%%%%%%%%%%%%%%%%%%%%%%%%%%%%%%%%%%%%%%%%%%%%%%%%%%%%%%%%%%%
\begin{document}
%%%%%%%%%%%%%%%%%%%%%%%%%%%%%%%%%%%%%%%%%%%%%%%%%%%%%%%%%%%%%%%%%%%%%%%

%%%%%%%%%%%%%%%%%%%%%%%%%%%%%%%%%%%%%%%%%%%%%%%%%%%%%%%%%%%%%%%%%%%%%%%
% Deja esto as\'{\i}: 
\previousvspace{-1.8cm}
\makesupertitle
\previousvspace{-0.25cm}
\maketitle
\vspace*{-0.4cm}
\thispagestyle{empty} % <<<<< Importante
%%%%%%%%%%%%%%%%%%%%%%%%%%%%%%%%%%%%%%%%%%%%%%%%%%%%%%%%%%%%%%%%%%%%%%%


Sea
\begin{equation*}
  f(x) = \frac{(x^2-x+1)^3}{x^2(x-1)^2}
\end{equation*}
y
\begin{align*}
  H_1(x) &= x, & 
  H_2(x) &= \frac{1}{x}, & 
  H_3(x) &= 1-x, \\ 
  H_4(x) &= \frac{1}{1-x}, & 
  H_5(x) &= \frac{x-1}{x}, & 
  H_6(x) &= \frac{x}{x-1}.
\end{align*}
Estas funciones cumplen $f(H_i(x)) = f(x)$, $i=1,\dots,6$, como vemos observando la recta punteada horizontal en el gr�fico:
\begin{center}
\begin{tikzpicture}[scale=1.25, y=0.19cm] % [scale=1.25,y=0.2cm] % [domain=-4:5]
    \draw[-to] (-4,0) -- (5,0) node[right] {$x$}; % Eje x
    \draw[-to] (0,0) -- (0,25) node[above] {$f(x)$}; % Eje y
%    \draw[thin, color=gray] (0,0) -- (0,23.1525); % As�ntota vertical en x = 0 (SOBRA)
    \draw[thin, color=gray] (1,0) -- (1,23.1525); % As�ntota vertical en x = 1
% La recta que pasa por los m�nimos: % (27/4 = 6.75)
    \draw[thin, color=gray] (-4,6.75) --
      (0,6.75) node[below left=-2pt, color=black] {$\frac{27}{4}$} -- (5,6.75);
% La recta que pasa a altura de f(3.75) = 13.612812:
    \draw[dotted] (-4,13.612812) 
      -- (5,13.612812) node[right, color=black]{$f(\lambda)$};
% Nota: Hay que dibujar antes estas cosas auxiliares para que 
% su color gris no "machaque" el negro de la gr�fica principal.
% La funci�n $f$ dibujada en tres trozos:
    \draw[thick, samples=50, domain=0.2:0.8] % centro
      plot(\x,{((\x)^2-\x+1)^3/((\x)^2*(\x-1)^2)}); 
    \draw[thick, samples=50, domain={-4}:{-0.25}] % izquierda 
      plot(\x,{((\x)^2-\x+1)^3/((\x)^2*(\x-1)^2)}); 
    \draw[thick, samples=50, domain=1.25:5] % derecha
      plot(\x,{((\x)^2-\x+1)^3/((\x)^2*(\x-1)^2)}); 
%% Marcas en el eje y a altura 27/4:
    \draw[shift={(0,6.75)}] (2pt,0pt) -- (-2pt,0pt);
%% Marcas sencillas en el eje x:
%  \foreach \x in {-4, -2, 0, 2, 4}
%    \draw[shift={(\x,0)}] (0pt,2pt) -- (0pt,-2pt) node[below] {$\x$}; 
%% Marcas m�s sofisticadas en el eje x:
    \foreach \x/\xtext in {-1/-1, 0.5/\frac12, 0/0, 1/1, 2/2}
      \draw[shift={(\x,0)}] (0pt,2pt) -- (0pt,-2pt) node[below] {$\xtext$}; 
%  \foreach \x/\xtext in {1/1, 1.5/1\frac{1}{2}, 2/2, 3/3}
%    \draw[shift={(\x,0)}] (0pt,2pt) -- (0pt,-2pt) node[below] {$\xtext$}; 
%% Marcas en el eje x correspondientes a $\lambda = 3.75$ y sus $H_i(\lambda)$
%% (en realidad, la parte "/xtext" no la uso pues no tengo "node", pero como 
%% es copiar y pegar de lo que sigue despu�s, da lo mismo):
    \foreach \x/\xtext in {-2.75/\scriptstyle 1-\lambda, 
        -0.36363636/\frac{1}{1-\lambda},
        0.26666667/\frac{1}{\lambda}, 
        0.73333333/\frac{\lambda-1}{\lambda}, 
        1.3636364/\frac{\lambda}{\lambda-1},
        3.75/\scriptstyle \lambda}
      \draw[shift={(\x,0)}] (0pt,2pt) -- (0pt,-2pt); 
%% L�neas punteadas verticales:
    \foreach \x/\xtext in {-2.75/\scriptstyle 1-\lambda, 
        -0.36363636/\frac{1}{1-\lambda},
        0.26666667/\frac{1}{\lambda}, 
        0.73333333/\frac{\lambda-1}{\lambda}, 
        1.3636364/\frac{\lambda}{\lambda-1},
        3.75/\scriptstyle \lambda}
      \draw[dotted, shift={(\x,0)}] 
        (0,13.612812) -- (0,-3.2) node[below=-2pt, color=black] {$\xtext$}; 
    \end{tikzpicture}
\end{center}
As� hemos mostrado que la funci�n $f(x)$ es invariante respecto al grupo de transformaciones $\{H_1(x),\dots,H_6(x)\}$. 
% (Seguro que es un resultado conocido, pero el que esto escribe no ha encontrado ninguna referencia.)

Se pueden construir tales funciones invariantes eligiendo cualquier polinomio de seis variables $Q(x_1,\dots,x_6)$ sim�trico y 
tomando $f(x) = Q(H_1(x),\dots,H_6(x))$. El ejemplo anterior se obtiene con $Q(x_1,\dots,x_6) = (x_1^2+\cdots+x_6^2+3)/2$.

%%%%%%%%%%%%%%%%%%%%%%%%%%%%%%%%%%%%%%%%%%%%%%%%%%%%%%%%%%%%%%%%%%%%%%%
\end{document}
%%%%%%%%%%%%%%%%%%%%%%%%%%%%%%%%%%%%%%%%%%%%%%%%%%%%%%%%%%%%%%%%%%%%%%%


